% =============================================================
% บทที่ 3 การติดตั้งระบบสำหรับพัฒนา
% ข้อมูลตรวจจาก main repo: README.md, CLAUDE.md, .env.example ของทั้ง 3 ส่วน, backend/scripts
% =============================================================

\section{ความต้องการของระบบ}

\subsection{ระบบปฏิบัติการ}

คู่มือนี้อธิบายการติดตั้งบนระบบปฏิบัติการ Windows 11

\subsection{ซอฟต์แวร์และบริการที่ต้องเตรียม}

\begin{parenlist}
\item Node.js เวอร์ชัน 20 ขึ้นไป (ระบบ CI ใช้เวอร์ชัน 20) สำหรับส่วนติดต่อผู้ใช้งานและบริการข้อมูลหลัก
\item Python เวอร์ชัน 3.11 สำหรับ AI Service
\item Git สำหรับดาวน์โหลดซอร์สโค้ดจาก GitHub
\item Visual Studio Code สำหรับเปิดและแก้ไขโครงงาน
\item บัญชี Supabase สำหรับฐานข้อมูล PostgreSQL
\item บัญชี Cloudinary สำหรับเก็บรูปภาพ
\item บัญชี Gmail ที่เปิดการยืนยันตัวตน 2 ขั้นตอนและสร้าง App Password สำหรับส่งอีเมล
\item โปรเจกต์ Google Cloud ที่เปิดใช้ Places API (New) และ Maps JavaScript API
\item คีย์ KKU AI Gateway สำหรับเรียกโมเดลภาษาขนาดใหญ่
\item คีย์ Serper API (ไม่บังคับ) สำหรับสคริปต์สร้างคำอธิบายสถานที่
\end{parenlist}

\section{ขั้นตอนการติดตั้งระบบสำหรับนักพัฒนา}

\subsection{การติดตั้ง Node.js}

\begin{numlist}
\item เข้าเว็บไซต์ https://nodejs.org/ แล้วเลือกดาวน์โหลดเวอร์ชัน LTS ดังภาพที่ \ref{fig:install-nodejs}
\item เปิดไฟล์ติดตั้งที่ดาวน์โหลดมา แล้วกด Next จนติดตั้งเสร็จสิ้น
\item เปิด PowerShell แล้วตรวจสอบการติดตั้งด้วยคำสั่งต่อไปนี้ ซึ่งควรแสดงหมายเลขเวอร์ชัน
\end{numlist}

\screenshot{install-nodejs}{หน้าดาวน์โหลด Node.js}

\begin{lstlisting}[numbers=none]
node -v
npm -v
\end{lstlisting}

\subsection{การติดตั้ง Python}

\begin{numlist}
\item เข้าเว็บไซต์ https://www.python.org/downloads/windows/ แล้วดาวน์โหลดตัวติดตั้ง Windows installer (64-bit) ของรุ่น 3.11 ขึ้นไป ดังภาพที่ \ref{fig:install-python}
\item ขณะติดตั้งให้เลือกตัวเลือก ``Add python.exe to PATH'' แล้วกด Install Now
\item ตรวจสอบการติดตั้งด้วยคำสั่ง \texttt{python --version}
\end{numlist}

\screenshot{install-python}{หน้าดาวน์โหลด Python}

\subsection{การติดตั้ง Git และ Visual Studio Code}

\begin{numlist}
\item ดาวน์โหลดและติดตั้ง Git จาก https://git-scm.com/downloads
\item ดาวน์โหลดและติดตั้ง Visual Studio Code จาก https://code.visualstudio.com/
\end{numlist}

\subsection{การดาวน์โหลดซอร์สโค้ดจาก GitHub}

เปิด PowerShell ไปยังโฟลเดอร์ที่ต้องการเก็บโครงงาน แล้วโคลน repository

\begin{lstlisting}[numbers=none]
git clone https://github.com/Poeir/dino-trip.git
cd dino-trip
\end{lstlisting}

โครงสร้างโฟลเดอร์ที่ได้ประกอบด้วย 3 ส่วนที่รันแยกจากกัน ดังตารางที่ \ref{tab:folders}

\begin{table}[H]
\caption{โครงสร้างโฟลเดอร์ของโครงงาน}
\label{tab:folders}
\centering
\begin{tabular}{|l|l|p{5cm}|}
\hline
\textbf{โฟลเดอร์} & \textbf{เทคโนโลยี} & \textbf{หน้าที่} \\
\hline
\texttt{frontend} & React 18 + Vite & หน้าเว็บที่ผู้ใช้และผู้ดูแลระบบใช้งาน \\
\hline
\texttt{backend} & Node.js + Express & REST API ฐานข้อมูล และสคริปต์นำเข้าข้อมูล \\
\hline
\texttt{chatbot-service} & Python + FastAPI & แชตบอต ระบบวางแผนทริป และการสร้าง Embedding \\
\hline
\end{tabular}
\end{table}

\subsection{การสร้างฐานข้อมูลบน Supabase}

\begin{numlist}
\item เข้าเว็บไซต์ https://supabase.com/ ลงชื่อเข้าใช้ แล้วสร้างโปรเจกต์ใหม่ พร้อมบันทึกรหัสผ่านฐานข้อมูลที่กำหนดไว้
\item เปิดหน้า Project Settings > Database > Connection string แล้วคัดลอก Session pooler (พอร์ต 5432) ไว้ใช้เป็นค่า \texttt{DATABASE\_URL} ไม่ควรใช้ Transaction mode (พอร์ต 6543) เพราะ backend เป็นเซิร์ฟเวอร์ที่รันต่อเนื่อง
\item เปิดหน้า Project Settings > API แล้วคัดลอก Project URL และ service\_role key ไว้ใช้กับ AI Service
\item สร้างโครงสร้างตารางด้วยการนำไฟล์ migration ไปใช้กับฐานข้อมูลใหม่ ดังคำสั่งในหัวข้อการนำเข้าข้อมูลตั้งต้น
\end{numlist}

\subsection{การติดตั้งไลบรารีของแต่ละส่วน}

เปิดโฟลเดอร์โครงงานด้วย Visual Studio Code (\texttt{code .}) แล้วใช้ Terminal ติดตั้งไลบรารีของแต่ละส่วนตามคำสั่งต่อไปนี้

\begin{lstlisting}[numbers=none]
cd backend
npm install

cd ../frontend
npm install

cd ../chatbot-service
python -m venv venv
venv\Scripts\activate
pip install -r requirements.txt
\end{lstlisting}

\subsection{การตั้งค่าไฟล์ .env}

แต่ละส่วนมีไฟล์ \texttt{.env.example} เป็นต้นแบบ ให้คัดลอกเป็นไฟล์ \texttt{.env} ในโฟลเดอร์เดียวกัน แล้วกรอกค่าตามตารางที่ \ref{tab:env-vars} ไฟล์ \texttt{.env} มีข้อมูลลับ จึงห้ามนำขึ้น Git

\begin{lstlisting}[numbers=none]
cd backend
copy .env.example .env
cd ../frontend
copy .env.example .env
cd ../chatbot-service
copy .env.example .env
\end{lstlisting}

\begin{table}[H]
\caption{ตัวแปรสภาพแวดล้อมที่ต้องกำหนดในไฟล์ .env}
\label{tab:env-vars}
\centering
\begin{tabular}{|l|l|p{5.2cm}|}
\hline
\textbf{ส่วนของระบบ} & \textbf{ตัวแปร} & \textbf{คำอธิบาย} \\
\hline
backend & \texttt{DATABASE\_URL} & connection string ของ Postgres (Session pooler พอร์ต 5432) \\
\hline
backend & \texttt{FRONTEND\_ORIGIN} & origin ของ frontend ต้องตรงกันทุกตัวอักษรและไม่มี / ท้าย เช่น \texttt{http://localhost:5173} \\
\hline
backend & \texttt{CHATBOT\_SERVICE\_URL} & ที่อยู่ของ AI Service เช่น \texttt{http://localhost:8000} \\
\hline
backend & \texttt{CLOUDINARY\_CLOUD\_NAME}, \texttt{CLOUDINARY\_API\_KEY}, \texttt{CLOUDINARY\_API\_SECRET} & ค่าจาก Cloudinary Dashboard > Settings > API Keys \\
\hline
backend & \texttt{GMAIL\_USER}, \texttt{GMAIL\_APP\_PASSWORD} & บัญชี Gmail และ App Password สำหรับส่งอีเมลยืนยัน \\
\hline
backend & \texttt{GOOGLE\_PLACES\_API\_KEY} & คีย์สำหรับนำเข้าและซิงก์สถานที่ (ต้องเปิด Places API (New) และผูก billing) \\
\hline
backend & \texttt{PORT} & พอร์ตของ backend ค่าเริ่มต้น 4000 \\
\hline
frontend & \texttt{VITE\_API\_URL} & ที่อยู่ของ backend เช่น \texttt{http://localhost:4000} \\
\hline
frontend & \texttt{VITE\_CHATBOT\_SERVICE\_URL} & ที่อยู่ของ AI Service เช่น \texttt{http://localhost:8000} \\
\hline
frontend & \texttt{VITE\_GOOGLE\_MAPS\_API\_KEY} & คีย์ Maps JavaScript API ควรจำกัดด้วย HTTP referrer \\
\hline
chatbot-service & \texttt{KKU\_API\_KEY} & คีย์ของ KKU AI Gateway \\
\hline
chatbot-service & \texttt{SUPABASE\_URL}, \texttt{SUPABASE\_SERVICE\_ROLE\_KEY} & ใช้ค้นข้อมูลสถานที่และค้นหาเชิงเวกเตอร์ \\
\hline
chatbot-service & \texttt{MODEL\_NAME}, \texttt{TRIP\_PLANNER\_MODEL\_NAME} & โมเดลของแชตบอตและระบบวางแผนทริป (ต้องตั้งบน production) \\
\hline
chatbot-service & \texttt{CORS\_ORIGINS} & origin ที่เรียก AI Service ได้ คั่นด้วย , และไม่มี / ท้าย \\
\hline
\end{tabular}
\end{table}

\subsection{การนำเข้าข้อมูลตั้งต้น}

ขั้นตอนนี้ทำเพียงครั้งเดียวตอนติดตั้งฐานข้อมูลใหม่ โดยใช้ Terminal ที่โฟลเดอร์ \texttt{backend}

\begin{lstlisting}[numbers=none]
cd backend
npx supabase db push --yes
npm run fetch:places
npm run import:places
npm run seed:events-knowledge
\end{lstlisting}

ความหมายของแต่ละคำสั่งมีดังนี้

\begin{numlist}
\item \texttt{npx supabase db push --yes} สร้างตารางและฟังก์ชันตามไฟล์ migration ใช้กับฐานข้อมูลใหม่เท่านั้น เพราะฐานข้อมูลที่ใช้งานจริงอยู่แล้วไม่มีประวัติ migration ทำให้คำสั่งนี้รัน migration เก่าซ้ำทั้งหมด
\item \texttt{npm run fetch:places} ดึงข้อมูลสถานที่จาก Google Places API ลงไฟล์ \texttt{backend/data/places.json} ซึ่งมีค่าใช้จ่ายตามจำนวนสถานที่
\item \texttt{npm run import:places} นำเข้าข้อมูลจากไฟล์ดังกล่าวเข้าฐานข้อมูล
\item \texttt{npm run seed:events-knowledge} เพิ่มข้อมูลกิจกรรมและฐานความรู้ตั้งต้น
\end{numlist}

หลังนำเข้าข้อมูลใหม่ทุกครั้ง ต้องสร้าง Embedding ใหม่ เพื่อให้แชตบอตและระบบวางแผนทริปค้นพบข้อมูล

\begin{lstlisting}[numbers=none]
cd chatbot-service
venv\Scripts\activate
python scripts/embed_content.py
\end{lstlisting}

\subsection{การสร้างบัญชีผู้ดูแลระบบ}
\label{sec:create-admin}

บัญชีผู้ดูแลระบบสร้างด้วยสคริปต์ที่โฟลเดอร์ \texttt{backend} โดยระบุอีเมลและรหัสผ่านที่ต้องการ

\begin{lstlisting}[numbers=none]
cd backend
node scripts/create-admin.js <email> <password>
\end{lstlisting}

เมื่อสำเร็จจะแสดงข้อความ ``Admin account ready'' และสามารถเข้าสู่ระบบที่หน้า \texttt{/admin/login} ได้

\subsection{การรันระบบ}

เปิด Terminal ใน Visual Studio Code แยก 3 หน้าต่าง แล้วรันแต่ละส่วนตามคำสั่งต่อไปนี้

\begin{lstlisting}[numbers=none]
# Terminal 1: Backend (http://localhost:4000)
cd backend
npm run dev

# Terminal 2: AI Service (http://localhost:8000)
cd chatbot-service
venv\Scripts\activate
uvicorn src.main:app --reload --port 8000

# Terminal 3: Frontend (http://localhost:5173)
cd frontend
npm run dev
\end{lstlisting}

เมื่อทุกส่วนทำงานแล้ว ให้เปิดเบราว์เซอร์ไปที่ \texttt{http://localhost:5173} จะปรากฏหน้าแรกของไดโน และสามารถดูเอกสาร API ของ AI Service ได้ที่ \texttt{http://localhost:8000/docs}

\subsection{การนำระบบขึ้นใช้งานจริง}

ระบบที่ใช้งานจริงที่ https://dinokhonkaen.app รันบนเซิร์ฟเวอร์เดียวด้วย Docker Compose ในโฟลเดอร์ \texttt{deploy} โดยมี Caddy ให้บริการ HTTPS เสิร์ฟหน้าเว็บ และส่งต่อเส้นทาง \texttt{/api} ไปยัง backend และ \texttt{/chat}, \texttt{/trip/llm} ไปยัง AI Service ทั้งนี้ การรวมโค้ดเข้ากิ่ง \texttt{main} บน GitHub จะทำให้ GitHub Actions เชื่อมต่อเซิร์ฟเวอร์และอัปเดตระบบอัตโนมัติ
